% ICIP/ICASSP (spconf) draft for MP-CycleGAN
% Build entry points:
%   - paper_blind.tex (double-blind)
%   - paper_camera_ready.tex (with author info)

\documentclass{article}
\usepackage{spconf,amsmath,graphicx}
\usepackage{booktabs}
\usepackage{multirow}
\usepackage{url}
\usepackage{cite}
\usepackage{enumitem}
\setlist[itemize]{leftmargin=*,noitemsep,topsep=2pt}

% --------------------
% Blind switch
% --------------------
\newif\ifICIPblind
\ifdefined\ICIPBLIND
  \ifnum\ICIPBLIND=1 \ICIPblindtrue \else \ICIPblindfalse \fi
\else
  \ICIPblindtrue
\fi

\title{MP-CycleGAN: Multi-Perceptual Driven Unpaired Underwater Image Enhancement}

\ifICIPblind
  \name{Anonymous Author(s)}
  \address{Anonymous Affiliation(s)}
\else
  % TODO: Replace with your real author list and affiliations
  \name{Author Name(s)}
  \address{Affiliation(s)}
\fi

\begin{document}
\maketitle

\begin{abstract}
Underwater images often suffer from color casts, haze, and contrast degradation due to wavelength-dependent absorption and scattering.
While supervised learning has achieved impressive performance, it heavily relies on paired training data that are difficult to acquire in real underwater environments.
This paper presents \emph{MP-CycleGAN}, a multi-perceptual driven unpaired enhancement framework built upon CycleGAN.
We incorporate (i) a gray-world based physical color constraint to mitigate wavelength-related color shifts, (ii) a gradient-based structure consistency loss to suppress structural distortions, and (iii) a VGG perceptual loss to preserve semantic content and textures.
On EUVP test samples (matched=200), MP-CycleGAN achieves 22.91 dB PSNR and 0.795 SSIM (95\% bootstrap CI reported), slightly outperforming a reproduced CycleGAN baseline.
We further evaluate cross-domain generalization from EUVP to UIEB (matched=890), demonstrating competitive perceptual quality under severe domain shift.
\end{abstract}

\begin{keywords}
Underwater image enhancement, unpaired learning, CycleGAN, physical prior, perceptual loss, domain generalization
\end{keywords}

\section{Introduction}
Underwater vision is a key prerequisite for oceanic engineering applications such as inspection, exploration, and autonomous navigation.
However, underwater images exhibit severe degradations including blue/green color casts, low contrast, and haze caused by wavelength-dependent attenuation and back-scattering.
These artifacts not only reduce visual quality but also hinder downstream tasks (e.g., detection and segmentation).

Deep models trained with paired degraded/clean images can yield strong results, yet collecting pixel-aligned pairs in real underwater scenes is expensive and often infeasible.
Unpaired learning (e.g., CycleGAN) provides an appealing alternative, but vanilla CycleGAN may introduce color drifting and structural artifacts when applied to underwater enhancement.
Moreover, full-reference metrics (PSNR/SSIM) and no-reference perceptual metrics (UCIQE/UIQM) often prefer different enhancement behaviors, leading to non-trivial trade-offs.

In this work we target the practical setting where only \emph{unpaired} data are available: underwater images and ``clean'' images are collected independently and are not pixel-aligned.
Compared with generic unpaired translation, underwater enhancement additionally requires stabilizing the global color statistics (to reduce wavelength-related casts) while preserving local structures under heavy scattering.

We propose \textbf{MP-CycleGAN}, a \textit{multi-perceptual} loss design that injects both physical prior and perceptual constraints into CycleGAN to improve stability and generalization.
Our main contributions are:
\begin{itemize}
  \item A multi-perceptual unpaired enhancement framework that combines a gray-world physical constraint, gradient-based structure consistency, and deep perceptual loss within CycleGAN.
  \item A unified evaluation on EUVP (in-domain) and UIEB (cross-domain), reporting 95\% bootstrap confidence intervals for PSNR/SSIM to improve result reproducibility.
  \item An ablation study analyzing the effect of each constraint under a controlled budget.
\end{itemize}

\section{Related Work}
\textbf{Physics/priors.} Classical methods correct underwater degradation by exploiting priors and imaging models, including wavelength compensation and dehazing~\cite{chiang2012wavelength}, fusion-based enhancement~\cite{ancuti2012fusion}, and restoration based on blurriness and light absorption~\cite{peng2017blurriness}.
Related haze-removal priors such as the dark channel prior (DCP)~\cite{he2009dcp} have also inspired underwater variants.
More recent work (e.g., Sea-thru) studies physically grounded color correction with calibrated range~\cite{akkaynak2019seathru}.
These approaches are interpretable but can be sensitive to water types and illumination changes.

\textbf{Supervised learning.} With paired datasets, CNN/GAN-based approaches can learn direct mappings and often achieve strong full-reference metrics.
UIEB~\cite{li2020uieb} provides paired references for 890 images and proposes WaterNet as a baseline.
FUnIE-GAN~\cite{islam2020funiegan} focuses on fast enhancement for underwater robotics, and UColor~\cite{li2021ucolor} improves color restoration via multi-color-space embedding.
However, supervision relies on limited paired data and may not generalize across different waters and cameras.

\textbf{Unpaired/weakly-supervised learning.} CycleGAN~\cite{zhu2017cyclegan} enables unpaired translation via adversarial learning and cycle consistency.
UGAN~\cite{fabbri2018ugan} and related methods explore adversarial learning for underwater enhancement with reduced supervision.
Nevertheless, standard unpaired translation may distort structures or produce unrealistic colors without additional constraints.

\section{Method}
\subsection{Overview}
Let domain $A$ denote underwater images and domain $B$ denote reference ``clean'' images.
CycleGAN learns two generators $G: A\rightarrow B$ and $F: B\rightarrow A$, and two discriminators $D_A, D_B$.
We keep the original adversarial and cycle-consistency objectives, and augment the enhancement direction with multi-perceptual constraints:
\begin{equation}
\mathcal{L} = \mathcal{L}_{GAN} + \lambda_{cyc}\mathcal{L}_{cyc} + \lambda_{idt}\mathcal{L}_{idt}
+ \lambda_{gray}\mathcal{L}_{gray} + \lambda_{struct}\mathcal{L}_{struct} + \lambda_{perc}\mathcal{L}_{perc}.
\end{equation}
Fig.~\ref{fig:arch} illustrates the framework.

\begin{figure}[t]
  \centering
  \includegraphics[width=\linewidth]{figures/mp_cyclegan_architecture.png}
  \caption{MP-CycleGAN overview. We enhance the CycleGAN framework with a gray-world physical constraint, a gradient-based structure loss, and a deep perceptual loss.}
  \label{fig:arch}
\end{figure}

\subsection{Gray-world physical loss}
Underwater images often exhibit dominant blue/green tones due to wavelength-dependent absorption.
Based on the gray-world assumption, we encourage the mean RGB channel responses of $G(x)$ to be balanced:
\begin{equation}
\mathcal{L}_{gray}(\hat{I})=\frac{1}{3}\sum_{c\in\{r,g,b\}}\left| \mu_c(\hat{I}) - \mu_{gray}(\hat{I}) \right|,
\quad
\mu_{gray}(\hat{I})=\frac{\mu_r(\hat{I})+\mu_g(\hat{I})+\mu_b(\hat{I})}{3}.
\end{equation}

\subsection{Structure consistency loss}
To reduce structural distortions during translation, we align image gradients between input $I$ and output $\hat{I}$:
\begin{equation}
\mathcal{L}_{struct}(I,\hat{I})=\left\|\nabla_x \hat{I}-\nabla_x I\right\|_1+\left\|\nabla_y \hat{I}-\nabla_y I\right\|_1.
\end{equation}
This term directly penalizes changes in edges and local textures.

\subsection{Perceptual loss}
We use a VGG-16 feature extractor $\phi(\cdot)$ to preserve high-level semantics:
\begin{equation}
\mathcal{L}_{perc}(I,\hat{I})=\left\|\phi(I)-\phi(\hat{I})\right\|_1.
\end{equation}
In our implementation, $\phi$ is taken from a fixed ImageNet-pretrained VGG-16 at an intermediate layer~\cite{johnson2016perceptual,simonyan2015vgg}.
Intuitively, $\mathcal{L}_{gray}$ stabilizes the global color tendency, $\mathcal{L}_{struct}$ discourages edge/texture drifting, and $\mathcal{L}_{perc}$ regularizes the translation in a semantic feature space.

\section{Experiments}
\subsection{Datasets and protocols}
\textbf{EUVP (in-domain).} We train on the EUVP unpaired split~\cite{islam2020funiegan} and report results on its test samples (matched=200).
\textbf{UIEB (cross-domain).} We evaluate generalization by directly testing on UIEB~\cite{li2020uieb} reference-890 (matched=890) without target-domain fine-tuning.

\subsection{Metrics}
We report full-reference PSNR/SSIM~\cite{wang2004ssim}, and no-reference UCIQE/UIQM~\cite{yang2015uciqe,panetta2016uiqm}.
For PSNR/SSIM, we additionally report 95\% bootstrap confidence intervals (CI) computed by the unified evaluation script (\texttt{bootstrap\_iters}=1000, \texttt{seed}=123).

\subsection{Implementation details}
We adopt a ResNet-9 generator and a PatchGAN discriminator (CycleGAN default setting).
All models use instance normalization and the standard CycleGAN image preprocessing (resize and random crop to $256{\times}256$).
CycleGAN is trained for 200 epochs with Adam ($\beta_1{=}0.5$), initial learning rate $2\times 10^{-4}$, and linear decay in the second half of training.
MP-CycleGAN is initialized from the CycleGAN checkpoint at epoch 200 and fine-tuned to epoch 202 with learning rate $1\times 10^{-4}$.
Loss weights are set to $\lambda_{gray}{=}0.1$, $\lambda_{struct}{=}2.0$, and $\lambda_{perc}{=}0.05$.
The ResNet-9 generator has 11.38M parameters, and runs at 16.99 ms/image for $256{\times}256$ input on a single CUDA device (batch size 1).

\subsection{Main results on EUVP}
Table~\ref{tab:euvp} reports the in-domain results on EUVP (matched=200).
MP-CycleGAN slightly improves PSNR/SSIM compared to CycleGAN and produces competitive perceptual quality.
We also note that non-reference perceptual metrics may respond differently to sharpening/contrast changes, so we report both UCIQE and UIQM.

\begin{table}[t]
\centering
\caption{In-domain evaluation on EUVP test samples (matched=200). PSNR/SSIM are reported as mean$\pm$half-width of 95\% bootstrap CI.}
\label{tab:euvp}
\small
\setlength{\tabcolsep}{4pt}
\begin{tabular}{lcccc}
\toprule
Method & PSNR $\uparrow$ & SSIM $\uparrow$ & UCIQE $\uparrow$ & UIQM $\uparrow$\\
\midrule
Input (Identity) & 20.33$\pm$0.39 & 0.794$\pm$0.008 & 24.88 & 5.77 \\
Gray-World & 19.42$\pm$0.43 & 0.784$\pm$0.009 & 24.27 & 6.72 \\
Gray-World+CLAHE & 16.59$\pm$0.25 & 0.702$\pm$0.010 & 27.24 & 6.06 \\
CLAHE & 16.93$\pm$0.25 & 0.697$\pm$0.008 & 28.01 & 6.14 \\
CycleGAN & 22.82$\pm$0.42 & 0.783$\pm$0.010 & 25.93 & 5.34 \\
MP-CycleGAN (Ours) & \textbf{22.91}$\pm$0.44 & \textbf{0.795}$\pm$0.011 & 25.82 & 5.41 \\
\bottomrule
\end{tabular}
\end{table}

\subsection{Ablation study}
We perform a controlled ablation by progressively adding each constraint.
Table~\ref{tab:ablation} shows that combining structural and perceptual constraints improves the SSIM/UIQM balance under limited training budget.

\begin{table}[t]
\centering
\caption{Ablation results (matched=50) using the benchmark summaries in your project.}
\label{tab:ablation}
\small
\setlength{\tabcolsep}{4pt}
\begin{tabular}{c l c c c c}
\toprule
ID & Losses & PSNR & SSIM & UCIQE & UIQM\\
\midrule
A & Baseline & 17.41 & 0.495 & 27.89 & 8.01 \\
B & +Gray & 17.63 & 0.495 & 28.82 & 8.19 \\
C & +Gray+Struct & 16.86 & 0.566 & 27.21 & 9.01 \\
D & +Gray+Struct+Perc & \textbf{18.30} & \textbf{0.584} & 28.44 & 9.34 \\
E & +Gray+Struct+Perc+Color & 16.95 & 0.572 & 27.63 & \textbf{9.42} \\
\bottomrule
\end{tabular}
\end{table}

\subsection{Cross-domain generalization on UIEB}
Table~\ref{tab:uieb} evaluates models trained on EUVP and tested on UIEB (matched=890).
We observe that MP-CycleGAN and CycleGAN exhibit different preferences between full-reference (PSNR/SSIM) and perceptual (UIQM) metrics under domain shift, consistent with typical unpaired enhancement behavior.
In particular, CycleGAN achieves slightly higher PSNR/SSIM, while MP-CycleGAN is marginally better on UIQM.

\begin{table}[t]
\centering
\caption{Cross-domain evaluation (EUVP$\rightarrow$UIEB, matched=890). PSNR/SSIM are mean$\pm$half-width of 95\% bootstrap CI.}
\label{tab:uieb}
\small
\setlength{\tabcolsep}{4pt}
\begin{tabular}{lcccc}
\toprule
Method & PSNR $\uparrow$ & SSIM $\uparrow$ & UCIQE $\uparrow$ & UIQM $\uparrow$\\
\midrule
Input & 17.36$\pm$0.29 & 0.773$\pm$0.008 & 21.70 & 6.80 \\
Gray-World & 16.95$\pm$0.23 & 0.774$\pm$0.008 & 21.78 & 7.75 \\
CLAHE & 17.65$\pm$0.23 & \textbf{0.821}$\pm$0.005 & 25.22 & 6.51 \\
Gray-World+CLAHE & 17.64$\pm$0.23 & 0.820$\pm$0.007 & 25.02 & 6.77 \\
Gamma & 16.82$\pm$0.25 & 0.769$\pm$0.008 & 20.51 & 6.95 \\
CycleGAN & \textbf{17.84}$\pm$0.19 & \textbf{0.626}$\pm$0.009 & \textbf{23.25} & 10.48 \\
MP-CycleGAN (Ours) & 17.54$\pm$0.19 & 0.624$\pm$0.009 & 22.46 & \textbf{10.50} \\
\bottomrule
\end{tabular}
\end{table}

\begin{figure}[t]
  \centering
  \includegraphics[width=\linewidth]{figures/visual_comparison_final_refined.png}
  \caption{Qualitative comparison on EUVP samples (Input / CycleGAN / MP-CycleGAN / GT).}
  \label{fig:qual}
\end{figure}

\begin{figure}[t]
  \centering
  \includegraphics[width=\linewidth]{figures/uieb_cross_domain_visuals.png}
  \caption{Qualitative comparison on UIEB under domain shift (trained on EUVP).}
  \label{fig:uiebqual}
\end{figure}

\subsection{Metric trade-off analysis}
Unpaired enhancement can exhibit an intrinsic trade-off between pixel-level fidelity (PSNR/SSIM) and perceptual sharpness/contrast (captured by UIQM components).
Fig.~\ref{fig:tradeoff} visualizes the observed PSNR--UIQM tendency across evaluated methods, highlighting that our constraints can improve structural similarity on EUVP while remaining competitive in UIQM under domain shift.

\begin{figure}[t]
  \centering
  \includegraphics[width=\linewidth]{figures/psnr_uiqm_tradeoff.png}
  \caption{PSNR--UIQM trade-off illustration based on the unified evaluation outputs.}
  \label{fig:tradeoff}
\end{figure}

\subsection{Reference to supervised baselines (not strictly comparable)}
To provide a broader context, Fig.~\ref{fig:sota} summarizes representative supervised baselines (WaterNet, FUnIE-GAN, UColor) reported in their original papers.
These results are \textit{not} directly comparable due to differences in training data, preprocessing, and evaluation protocols; our main focus is the unpaired setting and the fair comparison between the reproduced CycleGAN and MP-CycleGAN under the same scripts.

\begin{figure}[t]
  \centering
  \includegraphics[width=\linewidth]{figures/sota_comparison_refined.png}
  \caption{Reference performance of representative supervised baselines reported in the literature (for context only).}
  \label{fig:sota}
\end{figure}

\section{Conclusion}
We presented MP-CycleGAN, an unpaired underwater image enhancement method that integrates a gray-world physical prior and multi-perceptual constraints to mitigate color drifting and structural artifacts.
Experiments on EUVP and cross-domain UIEB evaluation indicate that MP-CycleGAN is competitive to a reproduced CycleGAN baseline and offers improved perceptual quality under domain shifts.
Future work includes stronger underwater-specific priors and task-driven evaluation with downstream perception models.

% NOTE: ICIP allows an optional 6th page for references only.
\bibliographystyle{IEEEtran}
\bibliography{refs}

\end{document}
